\chapter{Local cutoffs and interior graphical curvature}\label{sec:graphical-local-cutoffs}
We retain the flat hypersurface conventions, $n\geq1$, of Definition \ref{def:periodic-graph-flow}. In this section $\mathcal P=\partial_t^\perp-\Delta$, $v=\sqrt{1+|Du|^2}$ and $a=|A|^2$. The local estimates follow the method of \cite{EH91}, Theorem 2.1 and Theorem 3.1, printed pp.~550 and 554--557. Compact support of a cutoff is the ingredient replacing a hypothesis at infinity. The rational curvature weight below has an exact tensor-square identity before its negative terms are discarded.

\begin{lemma}[Spatial cutoffs on an entire graphical flow]\label{lem:graphical-compact-cutoff-geometry}
Let $F(x,t)=(x,u(x,t))$ be a smooth entire graph on a closed finite time slab $[0,T]$. For an ambient center $X_0$, set
\[
 \eta=R^2-|F-X_0|^2-2nt.
\]
Then $\mathcal P\eta=0$ and $|\nabla\eta|^2\leq4R^2$ wherever $\eta>0$. The set of base coordinates on which $\eta\geq0$, over all these times, lies in one compact base ball. For the base cutoff $\zeta=R^2-|x-y_0|^2$, one has instead
\[
 0\leq\mathcal P\zeta=2\operatorname{tr}(g^{-1})\leq2n,
 \qquad |\nabla\zeta|^2\leq4R^2\quad(\zeta>0).
\]
\end{lemma}
\begin{proof}
The Euclidean position equation gives $\mathcal P|F-X_0|^2=-2n$, proving the first identity. The intrinsic gradient is twice the tangential projection of $F-X_0$, giving its bound. An ambient ball bounds the horizontal coordinate by its projection onto the base, uniformly in time. On that compact base ball, smoothness bounds all relevant derivatives on the entire finite slab, even if they are unbounded globally. For the base cutoff, Lemma \ref{lem:graphical-heat-coordinate-cancellation} gives $\mathcal P|x-y_0|^2=-2\operatorname{tr}(g^{-1})$. The eigenvalues of $g^{-1}$ lie in $(0,1]$, and its quadratic form is at most the Euclidean one, proving both estimates. These are the two cutoffs used in \cite{EH91}, Corollary 3.2; no general sign assumption on the heat operator of an unspecified radius function is needed.
\end{proof}

\begin{theorem}[The localized reciprocal-angle estimate]\label{thm:graphical-local-angle-cutoff}
For a smooth entire graphical flow on $[0,T]$, every ambient center and every $R>0$ satisfy
\[
 v(x,t)\bigl(R^2-|F(x,t)-X_0|^2-2nt\bigr)_+
 \leq\sup_y v(y,0)\bigl(R^2-|F(y,0)-X_0|^2\bigr)_+.
\]
The right side is finite by compactness of the base support. This is \cite{EH91}, Theorem 2.1, p.~550, for an entire graph. No global curvature or slope bound is assumed.
\end{theorem}
\begin{proof}
On $\eta>0$ put $q=v\eta$. The angle equation and Lemma \ref{lem:graphical-compact-cutoff-geometry} give
\[
 \left(\mathcal P+2v^{-1}\nabla v\cdot\nabla\right)q=-a q\leq0.
\]
Indeed the two cross-gradient terms and the term $-2\eta v^{-1}|\nabla v|^2$ cancel after adding this drift. The function $v\eta_+$ is continuous and vanishes at the spatial boundary of its support. All of that support lies in a compact base set over a finite time interval. On its interior every coefficient is smooth and bounded. The graph-coordinate drift from Lemma \ref{lem:graphical-material-drift} vanishes at a spatial maximum, so the usual first-positive-maximum argument applies to $v\eta_+-M-\epsilon t$, where $M$ is its initial supremum. A positive maximum cannot occur on $\eta=0$. Let $\epsilon\downarrow0$. Only the positive region is differentiated; no smoothness of the positive-part function across its zero set is asserted.
\end{proof}

\begin{corollary}[Initial bounded slope propagates without a growth assumption]\label{cor:entire-slope-without-growth}
If a smooth entire graphical flow on $[0,T]$ has $v(\cdot,0)\leq V_0<\infty$, then $v\leq V_0$ on the slab. No polynomial-growth hypothesis on the later slope is necessary. This is \cite{EH91}, Remark 2.2.
\end{corollary}
\begin{proof}
The right side of Theorem \ref{thm:graphical-local-angle-cutoff} is at most $V_0R^2$. Fix a finite point and time and take $R$ large enough for the cutoff there to be positive. Dividing gives
$v(x,t)\leq V_0R^2/(R^2-|F(x,t)-X_0|^2-2nt)$. Let $R\to\infty$. The point is fixed during this limit, so its position is finite and the denominator divided by $R^2$ tends to one. This strengthens Proposition \ref{prop:entire-graph-slope-preservation}; its earlier sufficient growth hypothesis is not required by this localized argument.
\end{proof}

\begin{proposition}[A rational weight with a negative curvature square]\label{prop:graphical-rational-curvature-identity}
Put $s=v^2$, choose a constant $k>0$ with $ks<1$, and define
\[
 \phi(s)=\frac{s}{1-ks},\qquad f=\phi(s)|A|^2,
 \qquad b=\frac{\phi'(s)}{\phi(s)}\nabla s.
\]
Then the exact identity is
\[
 (\mathcal P+b\cdot\nabla)f
 =-2|\nabla(\sqrt{\phi}\,A)|^2-2kf^2
       -\frac{k|A|^2}{2(1-ks)^3}|\nabla s|^2.
\]
If $1\leq v\leq V$ and $k=(2V^2)^{-1}$, then $|b|\leq2\sqrt2\sqrt f$. The tensor formula is smooth also at $A=0$ and refines the inequality in \cite{EH91}, equation (20), p.~555.
\end{proposition}
\begin{proof}
The previous evolution formulas give
$\mathcal Pa=-2|\nabla A|^2+2a^2$ and
$\mathcal Ps=-2as-3|\nabla s|^2/(2s)$.
The product and chain rules, followed by adding the drift, give
\[
 (\mathcal P+b\cdot\nabla)f
 =-2\phi|\nabla A|^2+2a^2(\phi-s\phi')
 -a\left(\phi''+\frac{3\phi'}{2s}-\frac{(\phi')^2}{\phi}\right)|\nabla s|^2
 -\phi'\langle\nabla a,\nabla s\rangle.
\]
Expanding $-2|\nabla(\sqrt\phi A)|^2$ accounts for the first and last terms and for $-a(\phi')^2|\nabla s|^2/(2\phi)$. Direct differentiation gives
\[
 \phi-s\phi'=-k\phi^2,\qquad
 \phi''+\frac{3\phi'}{2s}-\frac{3(\phi')^2}{2\phi}
       =\frac{k}{2(1-ks)^3},
\]
proving the identity. For the drift estimate, Lemma \ref{lem:graphical-slope-gradient} gives $|\nabla s|^2\leq4s^2(s-1)a$. Hence
\[
 |b|^2\leq\frac{4(s-1)a}{(1-ks)^2}
       =\frac{4(s-1)}{s(1-ks)}f\leq8f,
\]
since $s\geq1$ and $ks\leq1/2$. This bound is independent of $V$ after measuring curvature by $f$.
\end{proof}

\begin{theorem}[Interior curvature in a fixed base ball]\label{thm:graphical-base-interior-curvature}
Suppose a smooth graphical flow is defined on a neighborhood of $\overline B_R(y_0)\times[0,T]$ and $1\leq v\leq V$ there. For $0\leq\theta<1$ and $0<t\leq T$,
\[
 \sup_{|x-y_0|\leq\theta R}|A(x,t)|^2
 \leq\frac{C(n)V^4}{(1-\theta^2)^2}
             \left(\frac1t+\frac1{R^2}\right).
\]
No bound on the initial curvature occurs in this estimate. This is \cite{EH91}, Theorem 3.1 and Corollary 3.2(ii), pp.~554--557.
\end{theorem}
\begin{proof}
Use $k=(2V^2)^{-1}$ and $f,b$ from Proposition \ref{prop:graphical-rational-curvature-identity}. Let $\zeta=R^2-|x-y_0|^2$, and maximize $Z=\tau\zeta^2 f$ over the closed base ball and $0\leq\tau\leq t$. It is zero at initial time and on the lateral boundary; all factors are bounded there by smoothness. If its maximum is positive, it occurs at an interior spatial point and a time $\tau>0$. There $\nabla f=-2f\nabla\zeta/\zeta$. Writing $D=\mathcal P+b\cdot\nabla$, the product rule at this point gives
\[
 0\leq DZ\leq\zeta^2 f-2k\tau\zeta^2f^2
       +2\tau\zeta f\mathcal P\zeta
       +6\tau f|\nabla\zeta|^2
       +2\tau\zeta f\langle b,\nabla\zeta\rangle.
\]
The left inequality remains true in graph coordinates since the total tangential drift vanishes at the maximum. Divide by $\tau f$. The drift bound and Young's inequality give
\[
 2\zeta|b||\nabla\zeta|
 \leq k\zeta^2f+\frac8k|\nabla\zeta|^2.
\]
Lemma \ref{lem:graphical-compact-cutoff-geometry} therefore yields
\[
 k\zeta^2 f\leq\frac{\zeta^2}{\tau}
                 +4nR^2+4(6+8/k)R^2.
\]
Multiplying by $\tau/k$ bounds the maximum of $Z$ by
$R^4/k+(4n+24+32/k)tR^2/k$. At the requested point $\zeta\geq(1-\theta^2)R^2$, and $a\leq f$ because $\phi\geq1$. Divide by $t\zeta^2$ and use $k^{-1}=2V^2$, $V\geq1$. This proves the asserted constant form. The argument is an interior estimate for an existing smooth flow, not a continuation theorem.
\end{proof}

\begin{corollary}[Curvature in a shrinking ambient ball]\label{cor:graphical-ambient-interior-curvature}
For a smooth entire graphical flow on $[0,T]$, suppose $v\leq V$ wherever $|F-X_0|^2+2n\tau\leq R^2$, $0\leq\tau\leq T$. On the smaller region $|F(x,t)-X_0|^2+2nt\leq\theta R^2$, $0\leq\theta<1$, $t>0$,
\[
 |A(x,t)|^2\leq\frac{C(n)V^4}{(1-\theta)^2t}.
\]
This is the ambient-cutoff form of \cite{EH91}, Corollary 3.2(i).
\end{corollary}
\begin{proof}
Repeat the proof of Theorem \ref{thm:graphical-base-interior-curvature} with $\eta=R^2-|F-X_0|^2-2n\tau$ and $Z=\tau\eta_+^2 f$. Lemma \ref{lem:graphical-compact-cutoff-geometry} supplies compact base support, $\mathcal P\eta=0$ and $|\nabla\eta|^2\leq4R^2$ in its positive region. Thus the same maximum argument applies, now without the $4nR^2$ term. At the target point $\eta\geq(1-\theta)R^2$, giving a bound by $C V^4(1-\theta)^{-2}(t^{-1}+R^{-2})$. If this region is nonempty, $2nt\leq\theta R^2<R^2$, so $R^{-2}\leq(2nt)^{-1}$. Absorb this term in $C(n)t^{-1}$. Values on the zero set of the cutoff are never differentiated.
\end{proof}
